% !TEX root = ../main.tex
\section{Methodology}
\label{sec:methodology}

\subsection{Dataset Preparation and Crop Extraction}
\label{sec:methodology-crops}

The original dataset provides one bounding box per image, marking the workpiece's
location and size. To study how classification behaviour changes as the object occupies
more or less of the frame, this thesis defines the \emph{object-to-image (O2I) ratio} as
the fraction of a cropped image's \emph{area} occupied by the object's bounding box, and
generates, for every image, a square crop centred on the bounding box centre such that the
object occupies a chosen target ratio of the crop's area. For a bounding box of pixel
width $w$ and height $h$, the side length of the crop achieving a target ratio $r$ is

\begin{equation}
\text{crop\_side} = \sqrt{\frac{w \times h}{r}}.
\label{eq:crop-side}
\end{equation}

Five target ratios were used: $90\%$, $70\%$, $50\%$, $30\%$, and $10\%$, evenly spaced
from almost-pure-object to object-dominated-by-background. This range and spacing was
chosen to give a readable ratio-vs-performance curve with enough points to distinguish a
gradual decline from a sudden cliff, while keeping the number of required training runs
manageable; a $10\%$ floor was used to avoid the object shrinking to only a handful of
pixels after resizing to the network's input resolution. Applying this procedure to all 15
materials at all 5 ratios yields $15 \times 5 = 75$ distinct sub-datasets, each internally
consistent (same camera, same background, same object, only the crop tightness differs).

\subsection{Experimental Design: One-vs-Rest Binary Classification}
\label{sec:methodology-ovr}

An initial design trained a single 15-class classifier per ratio (5 models total, one per
ratio, each choosing among all 15 materials). This thesis instead adopts a
\emph{one-vs-rest} design: for every (material, ratio) combination, a separate binary
classifier is trained to answer the narrower question ``is this crop material $X$, or not''
(target $=1$ for the material in question, target $=0$ for crops of any of the other 14
materials at that ratio). This produces $15 \times 5 = 75$ independent binary models in
total.

This reformulation was adopted because it lets each material's sensitivity to the O2I
ratio be read off directly and independently: with one detector per material, the
detectability of e.g.\ aluminium as the object shrinks can be tracked in isolation, rather
than being entangled with how the other 14 materials' decision boundaries happen to shift
in a shared 15-way classifier. Every one-vs-rest model uses the same fixed random seed
($42$) on every run, so that seed variation is never a confound when comparing across
materials or ratios -- the ratio (and, for a given run, the material) is the only thing
that changes between models.

\subsection{Training Configuration}
\label{sec:methodology-training}

All 75 models share an identical training configuration, summarised in
Table~\ref{tab:training-config}.

\begin{table}[h]
\centering
\begin{tabular}{lp{9cm}}
\hline
\textbf{Parameter} & \textbf{Value} \\
\hline
Input resolution & $256 \times 256$ \\
Epochs & 100 (primary); 200 for a supplementary check, Section~\ref{sec:results-200epoch} \\
Batch size & 16--32 \\
Optimiser & Adam \\
Random seed & 42 (fixed across every run) \\
Loss & Weighted binary cross-entropy (see below) \\
\hline
\end{tabular}
\caption{Training configuration used for every one-vs-rest model.}
\label{tab:training-config}
\end{table}

An earlier, smaller-scale comparison on the original 15-class design had found
$128\times128$ input crops to perform marginally better than $256\times256$ ($80.5\%$ vs.\
$79.5\%$ validation accuracy at 25 epochs) while training $2.4\times$ faster, and $128px$
was initially adopted on that basis. For the one-vs-rest redesign, the input resolution
was changed to $256\times256$; this required no architecture change, since the adaptive
average pooling layer described in Section~\ref{sec:background-cnn} makes
\texttt{MaterialCNN} input-size independent by construction.

\subsubsection{Why Accuracy Is Not a Meaningful Metric Here}
\label{sec:methodology-metrics}

Reformulating the problem as 75 binary one-vs-rest tasks has an important side effect:
every task now has a severe class imbalance, since each binary classifier sees roughly one
positive example (the target material) for every 14 negative examples (all other
materials), an imbalance of approximately $1{:}14$. Under this imbalance, plain accuracy
is not a meaningful metric: a trivial classifier that always predicts ``not the target
material,'' without having learned anything about the visual difference between materials,
already achieves roughly $93\%$ accuracy by construction. This was verified directly on
synthetic data reproducing the class balance: the always-negative baseline scored $93.3\%$
accuracy while its F1 score -- which specifically measures how well the positive
(target-material) class is identified -- was $0.000$.

To avoid this failure mode, three changes were made to the training and evaluation
procedure, all three of which are applied together on every run:
\begin{enumerate}
\item \textbf{Weighted loss.} The loss function upweights the minority (positive) class by
approximately $14\times$, roughly matching the class imbalance ratio, so that misclassifying
a true positive is penalised proportionally more than misclassifying a true negative.
\item \textbf{F1 and ROC-AUC as the reported metrics.} Accuracy is never reported for any
one-vs-rest result in this thesis; F1 score and ROC-AUC on the positive (target-material)
class are used throughout instead, since both remain informative under class imbalance.
\item \textbf{Best-epoch checkpointing by F1, not accuracy.} The model checkpoint saved for
each run is the epoch with the highest validation F1 score, rather than the epoch with the
highest validation accuracy, so that the saved model is the one that is actually best at
the task that matters.
\end{enumerate}

\subsubsection{Validation and Test Split Protocol}
\label{sec:methodology-split}

All model development, hyperparameter decisions, and epoch selection were performed using
a held-out validation split. The test split was set aside at the start of the project and
was not inspected or used for any decision-making during development; it is reserved
exclusively for the final, once-only reported evaluation, so that no reported number is
the product of repeated tuning against the same data it is evaluated on.

\subsection{Explainability Pipeline}
\label{sec:methodology-explainability}

Both occlusion sensitivity (Section~\ref{sec:background-occlusion}) and Grad-CAM
(Section~\ref{sec:background-gradcam}) were implemented against the trained one-vs-rest
models and applied consistently across all 75 (material, ratio) combinations.

\subsubsection{Occlusion Sensitivity: Implementation and a Methodological Correction}
\label{sec:methodology-occlusion-impl}

The implementation slides a $32\times32$ grey patch across each $256\times256$ input
image with a stride of 16 pixels, re-running the trained model at every patch position and
recording the change in the network's output for the target class relative to the
unoccluded image. This produces a coarse $15\times15$ grid of importance values, which is
then upsampled by interpolation to the full $256\times256$ resolution to form the final
heatmap.

This specific patch size and stride were not tuned empirically against alternatives, so it
is worth stating explicitly why they are a reasonable choice rather than an arbitrary one.
Four considerations motivate the $32\times32$/stride-16 configuration. First,
\emph{computational cost}: a stride of 16 over a $256\times256$ image yields
$\left(\frac{256-32}{16}+1\right)^2 = 15^2 = 225$ patch positions, and therefore 225
forward passes per image. A stride of 1 would require $225^2 \approx 50{,}625$ forward
passes per image -- more than two hundred times as many -- which is infeasible at the
scale used in this thesis (75 models $\times$ 10 sampled images per method,
Section~\ref{sec:methodology-centerfrac}). A stride of half the patch size (16 is half of
32) is a standard compromise: adjacent patches overlap by 50\%, which smooths the
resulting heatmap after upsampling, without the cost of a finer stride. Second,
\emph{scale relative to the network's own receptive field}: the CNN backbone
(Table~\ref{tab:cnn-architecture}) halves spatial resolution at each of its four
pooling stages ($256\to128\to64\to32\to16$), so a single unit in the final feature map
used by Grad-CAM already aggregates information from a neighbourhood on the order of
tens of input pixels. A $32\times32$ occlusion patch is comparable to this scale, meaning
a single occluded patch is capable of meaningfully perturbing the activations the network
itself treats as one local unit of evidence, rather than being small enough to disappear
into the receptive field of a single feature or large enough to blur together several
unrelated ones. Third, \emph{scale relative to the object across the studied O2I ratios}:
at the smallest ratio studied (10\%), the object occupies a square of side
$\sqrt{0.10}\times256\approx81$ pixels (Equation~\ref{eq:crop-side}), so a $32\times32$
patch covers under half of the object's width -- large enough to register a meaningful
occlusion, small enough that several patches are needed to cover the object fully, which
is what gives the heatmap genuine spatial resolution \emph{within} the object rather than
one patch swallowing it whole. At the largest ratio (90\%), the same patch is a small
fraction of the object, allowing the heatmap to resolve structure across both object and
background. Fourth, \emph{consistency with Grad-CAM}: the Grad-CAM heatmap
(Section~\ref{sec:methodology-gradcam-impl}) is itself produced at $32\times32$ resolution
before upsampling, so the two methods are compared at a matched spatial resolution rather
than one being systematically coarser than the other. Together these considerations were
judged sufficient to justify the configuration used throughout this thesis; they are not a
substitute for an empirical sweep over patch sizes, which remains a possible direction for
future work (Section~\ref{sec:discussion-future}) but was not carried out here.

An important implementation decision, discovered during initial testing rather than
assumed in advance, concerns exactly which network output is measured. The first
implementation measured the change in the model's \emph{softmax probability} for the
target class. On a confidently and correctly classified test image, this produced a
heatmap that was completely flat -- minimum, maximum, and mean all exactly $0.0$ -- even
though Grad-CAM, run on the same model and image, produced a clear and sensible heatmap.
The cause was diagnosed rather than assumed: softmax probability is bounded to the
interval $[0, 1]$, and for a sufficiently confident model, the probability for the correct
class rounds to exactly $1.0$ in 32-bit floating-point representation both before and
after occluding any single patch, so the measured difference is always $1.0 - 1.0 = 0.0$
regardless of how much the model's underlying evidence actually changes. The fix was to
measure the change in the raw \emph{logit} (the network's output before the softmax
nonlinearity is applied) instead of the probability. Because logits are unbounded, they do
not saturate this way; on the same model and image, the logit-based heatmap ranged from
$-2.405$ to $3.145$ with clear, interpretable structure (strongest on a reflective rim and
a highlight on the object itself). \textbf{Every occlusion sensitivity result reported in
this thesis (Section~\ref{sec:results}) is therefore computed on the raw logit, not the
softmax probability}, specifically to avoid this saturation failure mode on confidently
trained classifiers.

\subsubsection{Grad-CAM: Implementation}
\label{sec:methodology-gradcam-impl}

The Grad-CAM implementation attaches to the final $256$-channel convolutional feature map
identified in Table~\ref{tab:cnn-architecture}. For a given input image and target class,
one forward pass computes this feature map and the class score; one backward pass computes
the gradient of that score with respect to every unit of the feature map. The gradients
are globally average-pooled within each of the 256 channels to obtain one importance
weight per channel; the channels are combined into a single $32\times32$ map using these
weights, passed through a ReLU (retaining only positive contributions to the target
class), min-max normalised, and upsampled to $256\times256$ to match the input resolution.

\subsubsection{The \texorpdfstring{$\text{center\_frac}$}{center\_frac} Metric}
\label{sec:methodology-centerfrac}

To compare occlusion sensitivity and Grad-CAM quantitatively, rather than only visually,
each heatmap is reduced to a single number, $\text{center\_frac}$: the fraction of the
heatmap's total (min-max normalised) importance that falls inside a square box centred on
the image, sized so that the box's \emph{area} equals that image's O2I ratio,

\begin{equation}
\text{box\_side} = \sqrt{\frac{\text{ratio}}{100}} \times 256,
\label{eq:box-side}
\end{equation}

matching the same area-based ratio definition used for crop generation
(Equation~\ref{eq:crop-side}). A $\text{center\_frac}$ of $1.0$ means the heatmap's
importance is entirely concentrated inside the (ratio-sized) central box, i.e.\ fully
object-focused; a value of $0.0$ means it falls entirely outside, i.e.\ fully
background-focused. For the rare case of a completely flat heatmap (denominator
approximately zero, see Section~\ref{sec:results-robustness}), a neutral fallback value of
$0.5$ is used rather than an undefined ratio. For each of the 75 models, $\text{center\_frac}$
is computed for both methods on 10 sampled validation images per model (fixed seed 42, so
sampling is reproducible), giving $75 \times 10 = 750$ image analyses per method.
